\begin{abstract}
Whole-cell models simulate a living cell from its molecular parts. They fix every
parameter to one value, so each prediction is a single sample with no error bar.
The models cannot say how much of a prediction the data support. Inference is hard
because the models mix deterministic and stochastic formalisms, the stochastic paths
are latent, and no parameter is shared across the boundary between the two blocks.
We present \IC{}, a Julia framework that makes inference a
first-class operation in whole-cell modelling. Every parameter is a prior
distribution with a recorded source. Modules couple through typed edges, and the
framework checks the composition and lists every reduction from the published model.
The framework can run the same gene-expression interface as an exact stochastic
process or as a mean-field ODE, so the whole cell compiles to one differentiable ODE
for gradient-based inference, at a declared and measured cost. We test the framework on a reduced model of the minimal cell
\syn{}, with four ODE modules and three stochastic modules. An exact sampler passes
simulation-based calibration on a two-gene reference system over \pn{vseven.R}
replicate datasets. \placeholder{Approximate method result (B7): agreement with the
exact posterior on M0, cost ratio.} \placeholder{Full-scale result (B8): calibration
on the reduced syn3A model; metabolite data shrink a gene-expression parameter they
do not measure.} \placeholder{Uncertainty-budget sentence: x\% of the uncertainty in
a prediction comes from parameters that no measurement supports.}
\end{abstract}

\placeholder{Author summary (PLOS CB, $\le$ 200 words) or Highlights and eTOC blurb
(Cell Systems): write once the venue is chosen.}
